\chapter{处理伦理问题}
\chapterauthors{Bernard Lo、Deborah G. Grady}
\sourcepage{179}
涉及人的研究会引出伦理问题，因为人们为了推进科学知识、使他人受益，承受了不便与风险。公众既参加临床研究，又为之提供经费，需要相信研究遵循高水平的伦理标准。

本章先介绍研究监管的历史，再讨论指导涉及人的研究的伦理原则及美国联邦法规，尤其是机构审查委员会（IRB）批准与知情同意要求。最后讨论科研不端、作者资格、利益冲突，以及某些研究类型特有的问题。

\section{临床研究法规的历史}
现有法规和指南，是针对许多不当行为逐步发展起来的，包括第二次世界大战纳粹医师的“研究”、美国对囚犯、长期照护机构居民及其他易受伤害群体的研究，以及塔斯基吉研究（案例 7.1）。
\begin{bookbox}{案例 7.1\quad 塔斯基吉研究}
美国卫生与公众服务部于 1932 年开始塔斯基吉研究（Tuskegee study），以确定未经治疗梅毒的长期影响（1）。入组者是亚拉巴马州农村贫困、受教育少的黑人男性，获得食物、一些基本医疗及丧葬保险。研究者谎称正在为他们治疗梅毒，
\sourcepage{180}
例如把研究用腰椎穿刺说成“特别免费治疗”。二战期间抗梅毒抗生素可用，后来又被推荐为公共卫生措施，但研究者仍设法阻止参与者接受治疗。结果广泛发表，直至举报者长期努力才终止研究。1974 年，联邦政府因这一事件发布涉及人的研究法规，要求知情同意及 IRB 审查。1997 年，Bill Clinton 总统正式道歉。
\end{bookbox}

\section{伦理原则}
为指导涉及人的研究，人们提出了四项原则（2），塔斯基吉等研究都曾违反。

第一，\term{尊重个人}（respect for persons）：每个人都有权自行决定是否参加，包括知情同意，以及随时退出的自由。儿童、重度痴呆成人等不具决策能力者，如果代理人允许、本人不反对，也可参加。

第二，\term{有利原则}（beneficence）：研究获得的科学知识，必须超过参与者所承受的不便与风险，并尽量减少风险。风险不仅是干预的身体伤害，也包括泄密、污名及歧视等心理社会伤害。例如，可通过筛查潜在参与者、排除较可能受到伤害者，监测参与者有无不良作用，确保研究人员经过培训和认证，以及保护保密性，来降低参与研究的风险。有利原则有时难以评价，因为风险仅由少数参与者承受，潜在利益却可能及于大量患者甚至全社会。
\sourcepage{181}
I 期药物试验就是典型张力：常在少数健康者中检验新药是否损害肝、肾或血液等。参与者有助于发现有效治疗，却可能只有风险、没有获益机会。

第三，\term{公正原则}（justice）：研究负担与利益应公平分配。如果其他人群也能回答问题，就不应专门选低收入、教育有限、医疗获取不足或决策能力受损等处于不利地位、易受伤害的人。主要因为容易接触、合作和随访而研究他们，是不公平地利用他们。

公正也要求公平获得潜在利益。研究可能让人接触新疗法，因此不应受收入、保险或教育限制。儿童、女性及少数族裔在研究中代表不足，导致知识依据薄弱，也可能降低这些群体的医疗质量。NIH 资助的临床研究，除非有充分理由，必须让这些群体得到充分代表。

最后，塔斯基吉违反\term{如实告知}（truth-telling）原则：研究者谎报真正目的，并在已知青霉素有效后仍隐瞒治疗信息。如实告知不仅是不撒谎，还要求提供相关信息、解释潜在风险，并诚实说明参与的可能利弊。

\section{涉及人的研究的联邦法规}
联邦法规适用于所有联邦资助研究，以及将提交美国食品药品监督管理局（Food and Drug Administration，FDA）支持新药或器械申请的研究（3）。
\sourcepage{182}
多数大学还要求所属教职工的涉及人的研究遵守知情同意及 IRB 审查等基本规定，包括私人资助或校外研究。法规虽用“受试者”（subjects），我们更喜欢“参与者”（participants），因为这承认人能够参与研究，而不只是被实验或提供数据的对象。

几个重要定义如下：
\begin{itemize}
\item \textbf{研究}是“旨在发展或增进可推广知识的系统性探究”。以个别患者获益而非发表为目的、尚未证实的临床照护，不视为研究。有些质量改进项目可按研究处理，但多数符合下文豁免标准。
\item \textbf{人的受试者}是活着的个体，研究者通过与其互动取得资料，或取得可识别的私人信息。
\item \textbf{私人信息}是人们可合理期待未被观察、记录或公开的信息，例如病历。
\item 研究者能够或可能轻易确定本人身份的信息，称\textbf{可识别信息}。
\item 先前收集资料的身份对应表，若在分析前销毁，或研究者不能取得，则编码研究资料不视为可识别。
\end{itemize}

\subsection{机构审查委员会批准}
法规要求涉及人的研究由 IRB 批准，有时称独立伦理委员会、伦理审查委员会或研究伦理委员会。它评审方案，确保研究合乎伦理、参与者福祉和权利受保护，
\sourcepage{183}
有权批准、否决或要求修改。美国卫生与公众服务部下的人体研究保护办公室（Office for Human Research Protections）负责监督 IRB。大多数机构要求所有此类研究接受监督，即使没有联邦资助。委员虽多为研究者，也必须包括社区成员及熟悉研究法律伦理者。

批准时，IRB 必须确定（3）：
\begin{itemize}
\item 已尽量减少参与者风险；
\item 相对于预期利益及知识的重要性，风险合理；
\item 参与者选择公平；
\item 将从参与者或其法定授权代表取得知情同意；
\item 保密充分。
\end{itemize}
多数大学、医疗中心和医院有自己的 IRB，以本机构表格、程序和指南执行法规，不能向更高机构上诉。许多大学也与提供独立审查服务的商业 IRB 建立关系。

跨机构项目在美国各地点，依法应依赖单一 IRB 批准，有时称“负责审查的 IRB”或“中央 IRB”，以免重复审查延误合作研究。当地 IRB 应向中央提供州法规及利益冲突、同意书措辞等机构政策的信息。

IRB 和法规受到的批评包括，过度重视同意书而忽视同意过程，未充分审查设计及科学价值（4、5）。它们虽然应审查方案修改和监测不良事件，
\sourcepage{184}
通常却不核查实际实施是否遵守批准方案，许多还缺少资源和专长。因此，法规及 IRB 批准应视为最低伦理标准。最终，研究者的判断和品格，才是确保研究可接受最重要的因素。

\subsection{免于全体 IRB 审查的情形}
多数调查、访谈，以及既有去标识记录和样本的二次分析，可以豁免 IRB 审查（表 7.1）。理由是风险低、几乎所有人都会同意，而且逐人取得同意会使费用或难度过高。不过，许多 IRB 仍要求提交项目资料，确认符合豁免。
\begin{bookbox}{表 7.1\quad 可豁免联邦研究法规的研究}
\begin{enumerate}
\item 调查、访谈或公共行为观察，除非同时满足：参与者容易识别；且公开回答会带来法律责任，或损害声誉、财务状况或就业能力。例如药物成瘾、抑郁、人类免疫缺陷病毒（human immunodeficiency virus，HIV）感染高危行为或非法移民问题的问卷，不予豁免。
\item 既有记录、资料或样本研究，前提是资料已公开，例如州或联邦机构发布的数据；或者研究者记录方式使身份不易确定，例如无法取得资料与身份的编码对应表。
\item 正常教育实践的研究。
\end{enumerate}
\end{bookbox}
\sourcepage{185}
某些最小风险研究，也可由一位委员\term{快速审查}（expedited review），不必全体审查（表 7.2）。保护办公室网站列出适用类型（6）。法规中的\term{最小风险}（minimal risk），是日常生活或常规身体、心理检查通常遇到的风险；IRB 既需考虑严重程度，也需考虑发生概率。
\begin{bookbox}{表 7.2\quad 可接受 IRB 快速审查的研究}
\begin{enumerate}
\item 某些最小风险程序，包括：
\begin{itemize}
\item 静脉采血、唾液或痰采集、皮肤或黏膜拭子；
\item 常规临床非侵入性程序，如心电图、磁共振检查；使人接受辐射的 X 线检查需全体审查；
\item 使用临床目的收集的资料、记录或样本；
\item 未获豁免的调查或访谈。
\end{itemize}
\item 已批准方案的小改动。
\item 研究除数据分析，或取得参与者常规临床程序的随访资料外均已完成时，续期批准。
\end{enumerate}
\end{bookbox}

\subsection{知情且自愿的同意}
研究者必须取得参与者自愿的\term{知情同意}（informed consent）：向潜在参与者说明研究关键内容、参与意味着什么、潜在利弊及替代选择，再取得同意。
\subsubsection{向参与者披露信息}
\sourcepage{186}
法规要求涵盖若干方面：
\begin{itemize}
\item \textbf{研究性质。}明确告知这是研究、目的为何、由谁负责、招募何种人；不必说明具体假设。
\item \textbf{研究程序。}说明要求做什么、花多少时间、访视或联系频率，并指出哪些不属常规照护。随机化及盲法应以能够理解的语言解释。访谈和问卷应说明主题。
\item \textbf{潜在风险、利益及替代选择。}用日常语言描述医疗、心理社会及经济方面利弊，也说明不参加的其他选择，例如试验治疗是否能在研究外获得。
\end{itemize}
有人担心，告知常低估风险、夸大利益，也未说明试验目标（7）。新药或新程序有时被描述成能让参与者获益，但大多数新干预即使初步结果令人鼓舞，最终也未显著优于标准治疗。许多人高估个人获益（7），称为\term{治疗性误解}（therapeutic misconception）。应明确说明尚不知研究治疗是否更好，而且看似有希望的药也可能造成严重伤害。

\subsubsection{同意书}
\sourcepage{187}
通常需要签署并经见证的书面同意书，记录研究者与参与者讨论的知情同意过程已发生，并包含上述必要信息。也可使用简式同意书，声明必要内容已口头说明；此时必须有见证者在场并签字。IRB 通常提供希望研究者采用的模板，要求的信息可能多于联邦法规。

\subsubsection{参与者对所披露信息的理解}
参与者常严重误解研究目标及利益。讨论和书面材料应避免术语及复杂句式，采用更长时间讨论、简明短小且易懂的问答式同意书，并检查参与者是否理解关键特征等办法来促进理解（8）。

\subsubsection{同意的自愿性}
伦理上有效的同意，既知情又自愿。应尽量减少胁迫或不当影响，例如过高报酬、招募囚犯或研究者自己的学生。如果不当影响使人忽视风险或削弱拒绝能力，就有问题。拒绝不得损害医疗照护，参与者也必须可随时退出；这些保证须写入同意书。

\subsection{知情同意的例外}
\sourcepage{188}
有些科学上重要的研究，如果要求逐人知情同意，就很难甚至不可能进行。
\subsubsection{使用剩余去标识样本与资料}
\begin{bookbox}{案例 7.2\quad 新生儿血样研究}
出生不久，婴儿接受足跟采血，血滴在滤纸上筛查遗传病。多数州的强制筛查不需父母许可，因此样本覆盖全体新生儿。剩余样本对研究出生缺陷、早产的遗传原因、妊娠期环境暴露及基因环境相互作用很有价值（9）。
\end{bookbox}
用去标识样本研究，不要求知情同意及 IRB 审查（表 7.1），但许多 IRB 仍要求通知。许多期刊投稿时，也要求作者说明方案经 IRB 批准，或经其认定不需审查。

\subsubsection{豁免知情同意}
某些有价值项目需要可识别资料和样本，不能豁免 IRB，却可能豁免知情同意。
\begin{bookbox}{案例 7.2\quad 新生儿血样研究（续）}
团队希望用可识别新生儿血样测量母亲特定化学物暴露，研究与低出生体重、早产及围产期死亡的关系。可把样本与出生证明、死亡证明及医院记录关联（9）。
\sourcepage{189}
为有足够效能，需研究大量儿童，逐一取得家长或监护人同意并不可行。
\end{bookbox}
法规允许 IRB 在表 7.3 所有条件成立时豁免。例如，上述母亲环境暴露与低出生体重研究，多数 IRB 会豁免同意。
\begin{bookbox}{表 7.3\quad 可以豁免知情同意的研究}
\begin{enumerate}
\item 风险不超过最小风险，豁免或改变不会不利于参与者权利和福祉，且没有豁免就无法切实开展。
\item 涉及可识别信息或生物样本时，没有这些可识别材料就无法切实开展。
\item 适当时，研究后会向参与者补充相关信息。这允许某些含有欺骗的研究，例如同意时披露目的会损害研究效度。
\end{enumerate}
\end{bookbox}

\subsubsection{例外的理由}
某些重要研究风险极低，取得同意只是负担，不能多保护参与者。每位患者都曾受益于利用既有病历和样本的研究；从互惠的公平角度，得到利益者也应愿意参加类似的极低风险研究，以惠及他人。

\subsubsection{对例外的反对}
联邦法规虽允许未经家长许可使用去标识新生儿血样，公众却有明显反对。
\sourcepage{190}
\begin{bookbox}{案例 7.2\quad 新生儿血样研究（续）}
数州家长反对，未经同意或不给退出机会，就保存新生儿血样供未指明研究；其中两州提起诉讼。他们不反对筛查采血，而认为去标识并没有解决隐私及家长自主权受损的担忧。
\end{bookbox}
一些州现已要求，利用州筛查样本研究须取得同意。因此，联邦法规允许之事，在某些州，尤其敏感研究，可能不允许。

\subsubsection{缺乏决策能力的参与者}
不能知情同意时，幼儿应由父母或监护人许可，成人应由法定授权代表许可。方案还应额外审查，确认不能改在有同意能力的人中研究。

\subsection{尽量减少风险}
应预见并降低风险，例如排除特别容易发生不良事件者、监测事件、换用较少侵入性的测量。保密是重要一环。
\subsubsection{保密}
泄密可能造成污名或歧视，尤其涉及性态度和行为、饮酒用药、违法及精神疾病时。
\sourcepage{191}
个人标识保护不足还可能带来财务风险。应编码资料、保护或销毁身份对应表、限制访问身份信息的人员，并加强数据安全。

但不应作无条件保密承诺。审计、传票，或发现法律强制报告的儿童虐待、某些传染病及严重暴力威胁时，可能不得不披露。可以预见时，方案应规定现场人员如何处理，并告知参与者。

\term{保密证书}（certificates of confidentiality）提供法律保护，禁止向研究无关者披露可识别资料，包括因传票或法院命令披露，除非参与者同意（10）。联邦机构批准或资助时会自动签发；没有联邦资助也可向 NIH 申请：\url{https://grants.nih.gov/policy/humansubjects/coc/coc-nih-funded.htm}。

\subsubsection{HIPAA 健康隐私规则}
联邦健康隐私规则，通常以《健康保险可携带性与责任法》（Health Insurance Portability and Accountability Act）的缩写 HIPAA 称呼，保护日常医疗、计费及行政中收集的可识别信息，即受保护健康信息（protected health information，PHI）。用于研究须签署授权（11），在 IRB 知情同意书之外另行取得。例如，参加随机试验之外，把资料或样本收入未来研究资料库还须单独授权。不可识别资料及某些其他情形不需授权。
\sourcepage{192}
关于隐私规则及其与保护人的研究联邦法规的区别，应咨询 IRB。

\subsection{保护易受伤害的参与者}
\term{易受伤害的研究参与者}（vulnerable research participants）是在伦理上更易遭不当利用的人。法规把所有儿童和囚犯视为易受伤害，并通过限制高风险研究加以保护；但限制也削弱这些群体治疗有效性与安全性的资料。法规未提及者，也可能在特定项目中易受伤害。

按易受伤害原因建立的替代框架，有不少优点（12）。有些人难理解利弊，或受到不当影响与胁迫，因而不能保护自身利益。改善同意过程可以促进知情自愿决定，例如让患者权益代表、亲友解释研究性质、风险、利益及拒绝权。识别不同原因，便可采用针对性保护，让人们能够参加有助于改善医疗知识的研究。

患病者另有一种脆弱性：医疗不良事件风险较高。可通过干预先导研究、排除最高风险者及全程密切监测应对。

\subsubsection{儿童研究}
儿童不能自行知情同意，却承担风险，因此易受伤害；幼儿还更容易受某些伤害，例如影响神经发育的暴露。
\sourcepage{193}
必须取得家长许可，并在发育程度允许时，取得儿童\term{赞同}（assent）。法规 D 分部仅在可能直接使儿童受益，或风险只略高于最小风险且很可能带来对儿童疾病或状况至关重要、可推广的知识时，允许高于最小风险的研究。

对能形成自己意见的儿童，应把儿童赞同与家长许可分开寻求，提供能理解的信息，并留意使孩子难表达反对的家庭关系及文化期待（13）。

\subsubsection{孕妇研究}
孕妇干预研究同时涉及她和胎儿的利弊。孕妇本身不被定义为易受伤害者，但旨在保护胎儿的法规 B 分部也会影响她。一般只有研究旨在处理母亲或胎儿的健康风险，且胎儿风险最小时，才可纳入孕妇。若目的是使胎儿获益，除父亲无法联系、无行为能力，或妊娠由强奸、乱伦导致外，须取得孕妇及父亲同意（3）。

历史上孕妇被排除在大多数试验外，缩小了孕期医疗及药物胎儿影响的知识范围；法规还延续了孕妇不能为自己及胎儿作知情决定的误解。对于妊娠期尚无已知安全有效疗法的严重疾病，研究者及资助方应努力将孕妇纳入试验，
\sourcepage{194}
并向 IRB 提供可能胎儿异常的资料。

\subsubsection{囚犯研究}
囚犯可能不觉得能够拒绝，也可能受金钱、暂离监狱常规生活或假释考虑影响。法规 C 分部限制联邦资助研究类型，要求更严格 IRB 审查及卫生与公众服务部批准。许多机构把同样要求用于所有囚犯研究。近年来，为改善知识基础，倡导者要求将囚犯纳入 HIV、丙型肝炎等有希望新疗法试验。

\subsubsection{难以理解研究利弊者}
法规未覆盖的常见原因，包括决策能力受损、教育少、健康素养及数字理解能力不足。对阅读或记忆能力差者，可口头分小段说明，检查理解，并由亲属或权益代表鼓励提问、提供支持。

\subsubsection{权力差异}
养老院等机构居民可能觉得有参加压力，倾向服从控制日常生活者，不明白拒绝不会遭报复或危及其他生活方面。

研究者若同时是主治医师，患者可能担心拒绝会降低医师对她的关心。
\sourcepage{195}
可由他人取得同意，并确保医师不知道是否参加，以减轻问题。同样，学生及受训者可能被迫参与老师或上级的研究。IRB 可要求研究者招募其他人，且不得知道谁拒绝。

\subsubsection{社会与经济劣势}
低社会经济地位或难获医疗的人，可能为了报酬或照护接受在较富裕时本认为不可接受的风险。教育或健康素养低者，可能不理解信息，或更容易受别人影响。
\begin{bookbox}{案例 7.3\quad 低收入国家养老院居民研究}
COVID-19 大流行期间，研究者可能计划在一个长期合作、却无法取得疫苗的低收入国家，开展已获美国及其他国家紧急使用授权疫苗的随机安慰剂试验。当时，获授权疫苗在 75 岁以上养老院居民中的安全性与疗效尚未研究，而该群体感染及死亡风险很高。资助方是制造商，结果会直接惠及类似患者。
\end{bookbox}
当时，美国疾病控制与预防中心（Centers for Disease Control and Prevention，CDC）把专业护理机构居民列为最高接种优先级。批评者可能认为，
\sourcepage{196}
不向对照组提供公共卫生指南推荐的干预，是不合伦理的。
\begin{bookbox}{案例 7.3\quad 低收入国家养老院居民研究（续）}
资助方与主要研究者可能提议，在研究期间几乎无疫苗可用的低收入国家开展。若证实安全有效，研究结束后按 CDC 建议向安慰剂组提供有效疫苗。
\end{bookbox}
历史上，低收入国家曾作为试验场所，但即使干预证实有效安全，当地仍得不到。HIV 流行早期，一些新抗病毒药安慰剂试验在发展中国家开展；在已有有效抗逆转录病毒药可供对照组的发达国家，这样做并不合伦理。现在认为，若发展中国家患者研究结束后仍无法取得有效疗法，就不能在那里开展此类试验（14、15）。

今天，资源充足国家已有标准有效治疗时，在低收入国家使用安慰剂，只有在设计有充分理由、结果惠及当地，且研究后提供治疗有可信而非仅仅愿景式的计划时，才合乎伦理。

案例 7.3 的参与者有多方面脆弱性，可能妨碍知情自愿同意：
\begin{itemize}
\item 在低收入国家，研究以外很难获得疫苗，可能认为没有其他获益途径（16）。
\sourcepage{197}
\item 认知受损、教育少或健康素养低的居民，可能无能力知情同意。
\item 机构生活中自主权常受限制，例如谁能进房间、何时吃饭或清洁，使他们更易受到胁迫。
\end{itemize}
公正原则要求公平分配负担与利益。让低收入国家居民承担在研究者本国被视为不合伦理的研究风险和不便，并不公平，也可能被认为是剥削。

研究者及 IRB 可增加下列保护：
\begin{itemize}
\item 研究实施国，以及资助方和研究者所在国，都必须审查。应有养老院居民权益代表参加；美国 IRB 还应有熟悉当地文化和研究困难的人。
\item 同意披露必须可理解。这需要翻译为参与者语言、评估当地不熟悉的研究概念、由患者权益代表审阅，并在当地预试验，避免外国研究者未察觉的权力关系造成不当影响。IRB 可要求可信的当地社区组织派同意过程监督员旁听。
\item 除非东道国对养老院研究有充分保护，伦理上最好只招募能够自行知情、自愿同意者。
\item 资助方及东道国政府应事先约定，试验后公平提供证实有效的疫苗，例如免费提供给该国其他养老院居民，
\sourcepage{198}
并按制造成本提供额外剂量。谈判时，政府可咨询世界卫生组织、全球疫苗倡导组织、慈善机构及社区顾问委员会。
\end{itemize}

\section{研究者的责任}
\subsection{科研不端}
美国联邦科研诚信办公室（Office for Research Integrity）将\term{科研不端}（research misconduct）定义为捏造、篡改及剽窃（17）。
\begin{itemize}
\item \term{捏造}（fabrication）是编造并记录或报告结果。
\item \term{篡改}（falsification）是操弄资料、材料、设备或程序，或改变、删除数据及结果，使研究记录不能反映实际发现。
\item \term{剽窃}（plagiarism）是挪用他人想法、结果或文字而不适当注明来源。
\end{itemize}
按定义，不端必须是故意，知道行为错误，或鲁莽，知道有明显风险却置之不理；不包括诚实错误及合理科学分歧。联邦定义未涵盖重复发表、拒绝分享材料、粗疏方法及性骚扰等其他不当行为，机构应以其他政策处理。下例是严重伤害参与者的较近案例。
\begin{bookbox}{案例 7.4\quad 基于基因组的癌症疗效预测模型}
一位年轻研究者开发了几个以肿瘤基因表达预测化疗反应的检查，
\sourcepage{199}
发表于高影响期刊，并用于 NIH 资助试验。另一大学生物统计学家不能重复结果，发现敏感与耐药肿瘤标签颠倒、输入错误导致基因名单不对、部分测试样本重复使用，以及评估前未固定算法等严重问题。错误发表后，研究者仍未处理，NIH 也不能重复主要结果。

因此 NIH 终止三项使用模型的随机试验，最终撤回 14 篇论文。研究者被发现伪造简历后辞去教职；联邦办公室认定他在论文及联邦资助申请中篡改数据，构成不端（18）。
\end{bookbox}
其他广为人知的案例，包括故意捏造或更改资料，声称麻疹－腮腺炎－风疹疫苗与儿童自闭症有关（19），以及声称通过体细胞核移植获得人类干细胞系（20、21）。这些行为破坏公众和医师信任，并威胁公共资助。

出现指控时，联邦资助机构及研究者所在机构有责任公平及时调查（22）。举报者及被指控者的权利都须尊重：前者免受报复，后者应知晓指控并有回应机会。认定不端后的处罚，可包括暂停资助、取消未来申请资格、撤销职位，以及行政、学术、刑事或民事制裁。

\sourcepage{200}
\subsection{作者资格}
要有\term{作者资格}（authorship）（22），研究者必须同时：
\begin{itemize}
\item 对构思设计，或数据分析解释作出重要贡献；
\item 对起草或修改稿件作出重要贡献；
\item 在投稿前最终批准稿件。
\end{itemize}
\term{挂名作者}（guest authors）与\term{隐名作者}（ghost authors）都不合伦理。挂名或荣誉作者只作微不足道贡献，却列名于论文、摘要或其他出版物，例如只提供名气、接触参与者的机会、试剂、实验室协助或经费。在研究完成、数据分析及初稿写完后才加入成为作者，也不适当。

隐名作者则作出重要贡献，却不列名，常是药企或医学写作公司的雇员。隐去他们，会使读者低估公司作用。一项研究显示，高影响综合期刊原创研究中，25\% 有挂名作者，12\% 有隐名作者（23）。

谁应署名、顺序如何，常有争议，最好在动笔前明确讨论决定。工作职责调整时，应协商作者变更；已有如何得体协商的建议（24）。没有决定排序的统一标准，所以有的期刊会逐一描述作者贡献。

\subsection{利益冲突}
主要利益应是有效回答重要科学问题、保护参与者安全。其他利益如果可能使研究偏倚、损害客观性或公众信任，
\sourcepage{201}
就是\term{利益冲突}（conflicts of interest）（25、26）。
\subsubsection{类型}
\begin{itemize}
\item \textbf{经济利益冲突。}新药、器械及检查研究常由产业资助。经济安排可能使设计实施偏倚、过度解释阳性结果，或不发表阴性结果（25）。拥有专利或制造公司股权者，若疗法有效，可能获得巨大利益。咨询费、酬金或实物礼物，也可能使判断偏向公司产品。
\item \textbf{职业利益冲突。}专业声望及对某想法的智识投入等非经济激励，也可能使人偏向预想结果。
\end{itemize}

\subsubsection{应对相冲突的利益}
所有利益冲突都应披露；有些造成偏倚的可能性大到应管理或避免。
\begin{description}[leftmargin=2em,labelindent=0pt,style=nextline]
\item[减少偏倚。]良好试验中的常规预防措施有助约束：让研究者不知道个人接受何干预，避免结局评估偏倚；由没有冲突的独立数据与安全监察委员会（Data and Safety Monitoring Board，第 11 章）审查中期资料，发现明确利弊时终止；资助、摘要及论文同行评审也有帮助。
\item[分离冲突角色。]理想情况下，个人经济激励或智识承诺强烈者，不应担任规划、分析或解释的关键角色，应由无此关系者负责。
\sourcepage{202}
例如干预专利持有人不应任主要研究者。但现实中，药品及器械制造商常在其资助试验担任重要角色。
\item[掌握分析与发表。]药企资助研究中，担任第一或资深作者的学术研究者，应确保合同赋予原始数据与统计分析控制权，以及无论药物有效与否都可发表的自由（27）。研究者有伦理义务对全部研究负责。资助方可以审稿、建议，并在投稿前确保专利已申请，但不应有权否决、审查删改发表，或强迫特定措辞。
\item[披露。]NIH、其他资助及联邦机构、当地 IRB、会议和期刊，都要求提交资助申请、摘要或论文时披露。药企还须在 \url{https://openpaymentsdata.cms.gov/} 公开对美国个别医师的付款。单纯披露往往不足以处理严重冲突，但可遏制不当行为，并让编辑、评审及读者判断不当影响。
\item[管理。]冲突重大时，机构、资助方或 IRB 可要求额外保护，例如更密切监测同意过程，或改变相关研究者职责。
\item[禁止特定情形。]资助方或机构可禁止专利持有人或制造公司高管担任试验主要研究者。
\end{description}

\sourcepage{203}
\section{某些研究类型特有的伦理问题}
\subsection{随机临床试验}
随机试验虽是评价干预最有力的设计，却有两项特有伦理问题：研究者主动干预，以及谁接受哪种干预由偶然决定。随机分配的一个理由是存在\term{临床均衡}（equipoise），直观似清楚，实际难定义（28）。它不要求各组完全平衡，但何组更好应确有不确定性或争议，使参与者把照护交由随机化、而非个人医师决定时，不会明显受害。

干预的不良作用常未知，因此需仔细监测。多数试验应有独立数据与安全监察委员会，间歇审查；若干预造成统计及临床上显著伤害，有权建议停止（第 11 章）。

对照干预也有问题。已有有效标准照护时，对照组应获得；但短期、无严重风险的轻微自限性疼痛等研究，安慰剂可能合理。须告知研究外可用的有效干预。

如果已有令人信服的有效或有害结果，或因招募少、退出多、结局发生不足而无法回答问题，继续也不合伦理。委员会定期分析中期资料，可据此决定提前终止（29）。不应由研究者自行分析，因为揭盲后继续研究可能带来偏倚，
\sourcepage{204}
而研究者对继续或停止常有利益冲突。中期资料审查程序及统计停止规则，应在入组前规定。发展中国家的试验还有案例 7.3 那样的额外困境。

\subsection{既有样本与资料的研究}
血液及组织样本库让未来研究不必再采样。虽无身体风险，仍有伦理问题：无人能预见未来研究，因而为未指定用途取得同意有困难。即使不易识别个人，参与者也可能反对特定用途。若通过基因序列重新识别身份而泄密，可能造成污名或歧视。即使个人未受害，群体也可能受害，如案例 7.5。
\begin{bookbox}{案例 7.5\quad 使用为既往研究收集的样本}
亚利桑那大学研究者在糖尿病高风险的美国原住民 Havasupai 部落中，研究相关遗传标志物。虽然同意书写可研究行为或医学疾病，与部落及参与者沟通却集中在糖尿病。后来样本用于精神分裂症遗传及祖先迁徙模式研究。部落认为超出原同意，精神分裂症研究带来污名，进化遗传研究又与其起源叙事相抵触。部落索赔 5000 万美元，
\sourcepage{205}
最终以 70 万美元和解，大学道歉并归还剩余遗传样本（30）。
\end{bookbox}
现行规定下，如果不易确定身份，已收集的生物样本可不经同意研究。但有人建议，对临床不再需要的剩余血液、活检或手术样本，应预先询问患者是否同意二次研究（31）。

新采研究样本时，同意书应允许参与者同意或拒绝某些广泛未来类别。例如可选择：所有经 IRB 及科学评审批准的未来研究；仅特定疾病；或仅本次、不得未来使用。

也应告知可识别资料和样本是否共享给其他研究者，以及发现可能申请专利并开发成有经济价值的产品。

\subsection{人工智能与大数据研究}
\term{人工智能}（artificial intelligence，AI）使计算机执行通常由人完成的任务。\term{机器学习}（machine learning）是不再另外编程也能自动学习、改善表现的一类 AI。非临床例子包括人脸识别、推荐购买、筛选求职者及信用评估。它有望改善医疗及行政，例如诊断皮肤癌和眼病、解读影像。但过程像“黑箱”，难以理解，
\sourcepage{206}
若只依赖底层数据、未临床验证，可能产生错误、适用有限或有偏倚的方法。非医疗领域广用算法中，因建模和验证资料不具代表性而偏倚很常见。
\begin{bookbox}{案例 7.6\quad 预测高医疗费用患者的机器学习算法}
研究者以约 5 万人去标识电子病历，预测哪些人会有不良医疗结局，目的是向高风险者提供照护管理。她选择未来医疗费用作结局。成功的算法可改善高风险患者医疗，或通过减少再住院、急性期后照护等昂贵服务节约费用（32）。
\end{bookbox}
黑人患者除急诊及透析外，平均费用低于白人，很可能因医疗获取更少。按费用识别对象的算法，可能优先给白人服务，却忽略病情更重的黑人（32）。这让需要更大的人得不到有益照护，违反公正原则。若改用现患慢性病数量界定的健康状况，获选黑人比例会加倍，因为他们的慢病更多、更重；所识别患者占总费用、因而可能节约的比例也上升。因此，结局选择既可能造成偏倚，也可能降低效力。

开发和检验预测算法时，应使用能代表应用目标人群的训练及验证资料，
\sourcepage{207}
并确保训练集不包含既往有偏倚决策的病例。例如皮肤癌训练资料应包含表现常不同于白人的黑人及亚裔；冠状动脉病预测应避免因女性和黑人当年未被开具适当检查而有偏倚。还应在多样人群中测试并达到公平表现（33）。使用 AI 干预的临床试验方案已有共识指南（34）。

\subsection{传感器与移动健康设备试验}
智能手机应用及可穿戴传感器可以采集多种与健康有关的资料，包括心率和心律、血糖水平、身体活动，以及声音和言语模式。这些技术统称\term{移动健康}（mobile health），允许被动收集信息，已用于心律失常、高血压、糖尿病、帕金森病及抑郁等情况。移动健康可能有助于识别疾病危险因素、追踪自然史，并支持以移动健康指标为预测变量或结局变量的随机试验。但多数设备和应用不准确，或未接受准确度检验（35）；经常使用这些设备的人群也可能经过了高度选择。

下例展示低收入国家电子传感器研究的问题。虽然拟议研究涉及重大全球健康问题，仍有若干伦理困境。首先，监测影响隐私和保密，尤其同床者可能未同意、甚至不知道有研究。他们隐私受侵，泄密可能伤害他们，例如暴露婚外关系。当地文化对女性不忠的污名可能比男性更大。
\sourcepage{208}
若聘当地研究助理，保密尤其重要：这有助招募，并给研究村庄实际利益，却也增加助理认识参与者的可能。开展案例 7.7 这类研究前，应通过先导研究、社区顾问委员会及社区参与式研究，评估数字传感器、尤其视频的可接受性（36、37）。
\begin{bookbox}{案例 7.7\quad 用数字传感器研究疟疾}
一位疟疾研究者计划在疟疾高发的低收入国家，开展一项描述性研究，考察杀虫剂处理蚊帐的使用依从性（36）。这类蚊帐是标准公共卫生干预，但其实际使用情况研究不足。传感器比自报蚊帐使用情况能更准确地测量依从性。第一代传感器只能检测蚊帐是否展开；第二代传感器具有录像和接近传感功能，还能判断是否有人在蚊帐内，并识别是谁。了解哪些尚未感染、处于最高感染风险的儿童在蚊帐内，有助于设计更有效的预防干预。
\end{bookbox}

\section{其他问题}
\subsection{给参与者报酬}
临床研究参与者付出时间和精力，应得报酬，也应报销交通、照顾儿童等实际支出。实际中，为招募留存也可能需要补偿。对很不便或高风险的研究，常提高付款；但激励会引出不当诱导及公正问题。若高风险给更多钱，
\sourcepage{209}
社会经济地位低者可能违背更好的判断而冒险。为避免不当影响，有人建议只报实际费用，时间则按当地非技术劳动时薪补偿（38）。

\subsection{研究的临床影响}
职业早期可能为了建立学术地位而发表，但地位来自重要性而非单纯数量。若不回答能改善健康及生活质量的问题，却耗费资源、使参与者冒险；或采用薄弱设计、低质量实施，无法回答问题，都不合伦理。

\section{小结}
\begin{enumerate}
\item 研究者须确保遵守尊重个人、有利、公正及如实告知原则。
\item 须满足适用联邦规定，关键为知情同意及 IRB。应说明性质、风险、潜在利益及替代选择，并保密、遵守 HIPAA 隐私规则。
\item 儿童、囚犯、认知缺陷或社会不利者，须额外保护。
\item 研究者必须诚信，不得捏造、篡改或剽窃；披露并适当管理利益冲突；仅在作出重要智识贡献时署名，
\sourcepage{210}
也确保所有重要贡献者列为作者。
\item 特定研究另有问题：随机试验各组应存在均衡，对照干预适当，已证实某组更有效或有害时不得继续；既有样本资料尤其应保密；中低收入国家研究尤其应关注公正。
\end{enumerate}
\begin{references}
\item Jones JH, King NMP. Bad blood thirty years later: a Q\&A with James H. Jones. J Law Med Ethics. 2012;40(4):867-872.
\item The National Commission for the Protection of Human Subjects of Biomedical and Behavioral Research. The Belmont Report. Ethical Principles and Guidelines for the Protection of Human Subjects of Research [Internet]. HHS.gov. 1979 [cited 2021 Feb 12]. \url{https://www.hhs.gov/ohrp/regulations-and-policy/belmont-report/read-the-belmont-report/index.html}.
\item Department of Health and Human Services. Federal Policy for the Protection of Human Subjects. 45 CPR 46 [Internet]. Electronic Code of Federal Regulations (eCFR). 2017 [cited 2021 Feb 12]. \url{https://www.ecfr.gov/}.
\item Lo B, Barnes M. Federal research regulations for the 21st century. N Engl J Med. 2016;374(13):1205-1207.
\item Emanuel EJ, Menikoff J. Reforming the regulations governing research with human subjects. N Engl J Med. 2011;365:1145-1150.
\item Office of Human Research Protections. Expedited Review Procedures Guidance (2003) [Internet]. HHS.gov. 2003 [cited 2021 Feb 12]. \url{https://www.hhs.gov/ohrp/regulations-and-policy/guidance/guidance-on-expedited-review-procedures/index.html}.
\item Joffe S, Mack JW. Deliberation and the life cycle of informed consent. Hastings Cent Rep. 2014;44(1):33-35.
\item Nishimura A, Carey J, Erwin PJ, Tilburt JC, Murad MH, McCormick JB. Improving understanding in the research informed consent process: a systematic review of 54 interventions tested in randomized control trials. BMC Med Ethics. 2013;14(1):28.
\sourcepage{211}
\item Institute of Medicine. Challenges and Opportunities in Using Residual Newborn Screening Samples for Translational Research: Workshop Summary [Internet]. 2010 [cited 2021 Mar 4]. \url{https://www.nap.edu/catalog/12981/challenges-and-opportunities-in-using-residual-newborn-screening-samples-for-translational-research}.
\item Wolf LE, Beskow LM. New and improved? 21st century cures act revisions to certificates of confidentiality. Am J Law Med. 2018;44(2-3):343-358.
\item Nass SJ, Levit LA, Gostin LO, editors. Beyond the HIPAA Privacy Rule: Enhancing Privacy, Improving Health Through Research [Internet]. National Academies Press (US); 2009 [cited 2021 Feb 12]. \url{http://www.ncbi.nlm.nih.gov/books/NBK9578/}.
\item Gordon BG. Vulnerability in research: basic ethical concepts and general approach to review. Ochsner J. 2020;20(1):34-38.
\item Weisleder P. Helping them decide: a scoping review of interventions used to help minors understand the concept and process of assent. Front Pediatr. 2020;8:25.
\item HIV Prevention Trials Network. Updated Ethics Guidance for HIV Prevention Research [Internet]. 2020 [cited 2021 Jan 31]. \url{https://www.hptn.org/news-and-events/announcements/updated-ethics-guidance-hiv-prevention-research}.
\item The Ethics of Research in Developing Countries [Internet]. The Nuffield Council on Bioethics. 2002 [cited 2021 Jan 31]. \url{https://www.nuffieldbioethics.org/topics/research-ethics/research-in-developing-countries}.
\item Ndebele P. The Declaration of Helsinki, 50 years later. JAMA. 2013;310(20):2145-2146.
\item Office of Research Integrity. Definition of Research Misconduct | ORI—The Office of Research Integrity [Internet] [cited 2021 Feb 12]. \url{https://ori.hhs.gov/definition-misconduct}.
\item Institute of Medicine. Evolution of Translational Omics: Lessons Learned and the Path Forward [Internet]. 2012 [cited 2020 Dec 28]. \url{https://www.nap.edu/catalog/13297/evolution-of-translational-omics-lessons-learned-and-the-path-forward}. (See Appendix B).
\item Godlee F, Smith J, Marcovitch H. Wakefield’s article linking MMR vaccine and autism was fraudulent. BMJ. 2011;342:c7452.
\item Chong S, Normile D. How young Korean researchers helped unearth a scandal. Science. 2006;311(5757):22-25.
\item Chong S. Investigations document still more problems for stem cell researchers. Science. 2006;311(5762):754-755.
\sourcepage{212}
\item Mello MM, Brennan TA. Due process in investigations of research misconduct. N Engl J Med. 2003;349(13):1280-1286.
\item Wislar JS, Flanagin A, Fontanarosa PB, Deangelis CD. Honorary and ghost authorship in high impact biomedical journals: a cross sectional survey. BMJ. 2011;343:d6128.
\item Browner WS. Authorship. In: Publishing and Presenting Clinical Research. 3rd ed. Wolters Kluwer; 2013.
\item Lo B, Field M, editors. Conflict of Interest in Medical Research, Education, and Practice [Internet]. 2009 [cited 2021 Feb 13]. \url{https://www.nap.edu/catalog/12598/conflict-of-interest-in-medical-research-education-and-practice}.
\item Fineberg HV. Conflict of interest: why does it matter? JAMA. 2017;317(17):1717-1718.
\item DeAngelis CD, Fontanarosa PB. Ensuring integrity in industry-sponsored research: primum non nocere, revisited. JAMA. 2010;303(12):1196-1198.
\item Joffe S, Miller FG. Equipoise: asking the right questions for clinical trial design. Nat Rev Clin Oncol. 2012;9(4):230-235.
\item Ellenberg SS, Fleming TR, DeMets DL. Data Monitoring Committees in Clinical Trials: A Practical Perspective. 2nd ed. Wiley; 2019:496.
\item Mello MM, Wolf LE. The Havasupai Indian tribe case—lessons for research involving stored biologic samples. N Engl J Med. 2010;363(3):204-207.
\item Wolinetz CD, Collins FS. Recognition of research participants’ need for autonomy: remembering the legacy of Henrietta Lacks. JAMA [Internet]. 2020 [cited 2020 Nov 6];324:1027-1028. \url{https://jamanetwork.com/journals/jama/fullarticle/2769506}.
\item Obermeyer Z, Powers B, Vogeli C, Mullainathan S. Dissecting racial bias in an algorithm used to manage the health of populations. Science. 2019;366(6464):447-453.
\item Zou J, Schiebinger L. AI can be sexist and racist—it’s time to make it fair. Nature. 2018;559(7714):324-326.
\item Rivera SC, Liu X, Chan A-W, et al. Guidelines for clinical trial protocols for interventions involving artificial intelligence: the SPIRIT-AI Extension. BMJ [Internet]. 2020 [cited 2020 Dec 27];370. \url{https://www.ncbi.nlm.nih.gov/pmc/articles/PMC7490785/}.
\item Freeman K, Dinnes J, Chuchu N, et al. Algorithm based smartphone apps to assess risk of skin cancer in adults: systematic review of diagnostic accuracy studies. BMJ [Internet]. 2020 [cited 2021 Feb 13];368. \url{https://www.ncbi.nlm.nih.gov/pmc/articles/PMC7190019/}.
\sourcepage{213}
\item Krezanoski P, Haberer J. Objective monitoring of mosquito bednet usage and the ethical challenge of respecting study bystanders’ privacy. Clin Trials Lond Engl. 2019;16(5):466-468.
\item Fairchild AL. Objective monitoring of mosquito bednet usage and the ethical challenge of privacy revelations about study bystanders: ethical analysis. Clin Trials Lond Engl. 2019;16(5):469-472.
\item Gelinas L, Largent EA, Cohen IG, Kornetsky S, Bierer BE, Lynch HF. A framework for ethical payment to research participants. N Engl J Med [Internet]. 2018 [cited 2021 Feb 13]. \url{https://www.nejm.org/doi/10.1056/NEJMsb1710591}.
\end{references}

\chapterappendix{7A}{第 7 章习题：处理伦理问题}
\sourcepage{214}
\begin{enumerate}
\item 研究者想识别与 II 型糖尿病风险增加相关的基因。已完成的冠状动脉病危险因素前瞻性队列，保存冷冻血样及临床资料，包括研究开始时饮食、运动、临床特征、胆固醇及糖化血红蛋白，以及冠心病结局和新发糖尿病随访。拟进行 DNA 测序，不需新采血。
\begin{enumerate}
\item 可否依原队列知情同意开展？
\item 若原同意未许可，该怎样开展？
\item 设计采血研究时，怎样预先安排未来利用资料和样本？
\end{enumerate}
\item 某新药治疗晚期结肠癌初显前景，研究者计划 III 期随机试验。为减少人数，她希望用安慰剂，而非与现行治疗比较。
\begin{enumerate}
\item 安慰剂对照有何伦理问题？
\item 能否以合乎伦理的方式进行安慰剂研究？
\end{enumerate}
\item 为未来 HIV 疫苗试验作准备，研究者拟检验：接受最先进预防咨询后，能否招募仍有高血清转换率的队列；随访率是否足以开展疫苗试验。对象为注射吸毒者、以性交易换钱者、多性伴侣者等高风险人群，多数阅读及健康素养低。观察性队列随访两年，确定血清转换及随访率。
\begin{enumerate}
\item 法规要求知情同意披露什么？
\item 如何确保真正知情？
\sourcepage{215}
\item 观察期间，研究者有何责任降低高风险参与者 HIV 感染风险？
\end{enumerate}
\end{enumerate}
